\chapter{Cuadernos computacionales}

La comunidad científica fue de los primeros grupos en adoptar el uso profesional y a
gran escala de computadoras electrónicas.
Desde simulaciones numéricas de miles de cuerpos, análisis estadístico de miles
o millones de datos e incluso cómputo algebraico, dudo que exista un área de investigación
en la que el cómputo no sea importante en el desarrollo de nuevos resultados.
Por lo tanto, toda persona dedicada a la ciencia haría bien en familiarizarse con
el uso correcto de computadoras y con cómputo científico básico.

Hoy en día, el cómputo científico se encuentra en una intersección particular, a veces
desafortunada.
Para escribir código científico se requiere tanto \emph{conocimiento de desarrollo de software}
como \emph{conocimiento de dominio científico}, para implementar y validar lo que se desea
obtener.
Una persona típica que desarrolla \software~es experta en escribir código
utilizable, eficiente y mantenible.
Una persona que se dedica a la ciencia es experta en su dominio, y puede describir
con formalidad los objetos y modelos de su interés.
El conocimiento de ambas tiene \textit{a priori} una intersección vacía.\footnote{
  Ver por ejemplo\\
  \url{https://youtu.be/gWddFxOXefo?si=_hL469iAiF7LtahM} y\\
  \url{https://youtu.be/6m6ZkafwGDs?si=KBR94UFrty_E30EY}\\
  para una discusión informal en el contexto de resultados matemáticos.
}
Tal vez la persona científica esté familiarizada con herramientas computacionales de uso común
en su especialidad, pero hay una gran diferencia entre código que sirve para uso personal
y código apto para ser escrutinizado y reutilizado por otras personas (no se hable sobre
poder reproducir los resultados de otras personas).

Pensemos además en las diferencias entre la forma tradicional de presentar un resultado
científico y de presentar un proyecto de \software. Muchas veces, el estado del software
no es tan interesante para la comunidad---importan más los resultados obtenidos, su análisis
y su interpretación. Al mismo tiempo, como cualquier parte de la metodología de una
investigación, el código fuente de un proyecto debe ser escrutinizable.
Considerando todo esto, personalmente no me gustaría tener que revisar un archivo
\texttt{.c}, \texttt{.py} o \texttt{.f90} de \num{500} líneas de código y \num{2000} líneas
de comentarios explicando cómo funciona el programa y cómo usarlo.

En contextos académicos, las formas por excelencia de presentar trabajo,
proyectos y resultados son el artículo y la tarea. En ambas, se debe encontrar un
equilibrio entre ser demasiado conciso y demasiado verboso.
En un artículo se motiva la investigación, se presenta el contexto en que se sitúa el
trabajo, se explica la metodología y se discuten los resultados.
En una tarea usualmente pueden presentarse únicamente las respuestas o los resultados,
pero estos deben estar acompañados de prosa que ilustre qué se necesita para encontrar
la solución y de evidencia de que se realizó y entendió el trabajo.

Sería ideal, entonces, un formato que permitiera combinar la claridad de la prosa
explicativa con la ejecución del código relevante, en la misma forma en que se combinan
prosa y ecuaciones en documentos científicos. Este es exactamente el formato del
cuaderno científico.

\section{Cuadernos en general}


\section{Jupyter en particular}


\subsection{Escritura de cuadernos}

